\chapter{คอมพิวเตอร์วิทัศน์ตรวจวินิจฉัยโรคกุ้งและการจำแนกสเปกตรัมน้ำเขียว}
\label{chap:ch08}

\section*{แผนบริหารการสอนประจำบทที่ 8}
\addcontentsline{toc}{section}{แผนบริหารการสอนประจำบทที่ 8}

\begin{planbox}{หัวข้อเนื้อหาประจำบท}
\begin{enumerate}[topsep=2pt, itemsep=1pt, partopsep=0pt, parsep=0pt, leftmargin=1.5em]
    \item ระบาดวิทยาและพยาธิวิทยาของโรคสำคัญในกุ้งขาวแวนนาไมภาคตะวันออก
    \item สัณฐานวิทยาของตับ ตับอ่อน และทางเดินอาหารของกุ้งในสภาวะติดเชื้อ
    \item การพัฒนาโมเดลโครงข่ายประสาทเทียมน้ำหนักเบาเพื่อจำแนกโรคกุ้งบนสมาร์ตโฟน
    \item การตรวจติดตามสเปกตรัมการดูดกลืนแสงของแพลงก์ตอนพืชและปรากฏการณ์น้ำเขียว
    \item ระบบพยากรณ์ล่วงหน้าเพื่อป้องกันภาวะแพลงก์ตอนดรอปและสารพิษไซยาโนแบคทีเรีย
\end{enumerate}
\end{planbox}

\begin{planbox}{วัตถุประสงค์เชิงพฤติกรรม}
\begin{enumerate}[topsep=2pt, itemsep=1pt, partopsep=0pt, parsep=0pt, leftmargin=1.5em]
    \item ระบุลักษณะทางพยาธิสภาพของโรคตัวแดงดวงขาวและโรคตายด่วนในกุ้งขาวได้อย่างถูกต้อง
    \item พัฒนาสถาปัตยกรรมโครงข่าย MobileNet เพื่อจำแนกภาพพยาธิสภาพของตับกุ้งได้
    \item อธิบายความสัมพันธ์ของสเปกตรัมแสงกับการประเมินความเข้มข้นของคลอโรฟิลล์เอได้
    \item ออกแบบระบบเฝ้าระวังและเตือนภัยก่อนเกิดปรากฏการณ์น้ำล้มได้อย่างมีประสิทธิภาพ
\end{enumerate}
\end{planbox}

\newpage

\begin{planbox}{วิธีสอนและกิจกรรมการเรียนการสอน}
\begin{enumerate}[topsep=2pt, itemsep=1pt, partopsep=0pt, parsep=0pt, leftmargin=1.5em]
    \item บรรยายชีววิทยาโรคสัตว์น้ำและหลักการการตรวจวินิจฉัยด้วยภาพถ่ายดิจิทัล
    \item สาธิตการถ่ายภาพซูมระดับมาโครบริเวณตับและลำไส้กุ้งด้วยกล้องติดเลนส์ขยาย
    \item ปฏิบัติการฝึกสอนโมเดลจำแนกภาพพยาธิสภาพกุ้ง 4 คลาสด้วย Transfer Learning
    \item การวิเคราะห์ภาพถ่ายสเปกตรัมสีน้ำในบ่อเลี้ยงกุ้งเพื่อจำแนกชนิดของแพลงก์ตอน
\end{enumerate}
\end{planbox}

\begin{planbox}{สื่อการเรียนการสอน}
\begin{enumerate}[topsep=2pt, itemsep=1pt, partopsep=0pt, parsep=0pt, leftmargin=1.5em]
    \item เอกสารคำสอนบทที่ 8 สไลด์บรรยาย และคลังภาพถ่ายพยาธิสภาพโรคกุ้งมาตรฐาน
    \item กล้องดิจิทัลมาโครและชุดเลนส์ขยาย 20 เท่าสำหรับติดสมาร์ตโฟน
    \item สเปกโตรโฟโตมิเตอร์ชนิดพกพาและชุดฟิลเตอร์กรองแสงหลายย่านความถี่
    \item คอมพิวเตอร์ประมวลผลโมเดล Deep Learning (PyTorch และ TensorFlow Lite)
\end{enumerate}
\end{planbox}

\begin{planbox}{การวัดและประเมินผล}
\begin{enumerate}[topsep=2pt, itemsep=1pt, partopsep=0pt, parsep=0pt, leftmargin=1.5em]
    \item การทดสอบความรู้เรื่องโรคสัตว์น้ำและอาการทางพยาธิสภาพ
    \item การประเมินค่าความแม่นยำ (Accuracy และ F1-Score) ของโมเดลจำแนกโรคกุ้ง
    \item การประเมินความถูกต้องของแบบจำลองการทำนายปรากฏการณ์แพลงก์ตอนบลูม
\end{enumerate}
\end{planbox}

\clearpage

\section{ระบาดวิทยาและโรคสำคัญในกุ้งขาวแวนนาไม}

การระบาดของโรคติดเชื้อในฟาร์มเลี้ยงกุ้งแถบชายฝั่งทะเลภาคตะวันออกสร้างความเสียหายทางเศรษฐกิจอย่างมหาศาล โรคสำคัญที่เกษตรกรต้องเฝ้าระวังอย่างใกล้ชิดประกอบด้วย
\begin{enumerate}[leftmargin=1.5em]
    \item \textbf{โรคตัวแดงดวงขาว (White Spot Syndrome Virus หรือ WSSV)\index{โรคตัวแดงดวงขาว (WSSV)}\index{White Spot Syndrome Virus (WSSV)}:} เกิดจากเชื้อไวรัสดีเอ็นเอสายคู่ขนาดใหญ่ กุ้งที่ติดเชื้อจะมีจุดสีขาวขนาดเส้นผ่านศูนย์กลาง $0.5 - 2.0\text{ mm}$ ใต้เปลือกส่วนหัวและลำตัว ลำตัวเปลี่ยนเป็นสีชมพูแดง อัตราการตายสูงถึงร้อยละ 100 ภายใน $3 - 5\text{ วัน}$
    \item \textbf{โรคตายด่วนหรือตับวายเฉียบพลัน (EMS / AHPND\index{โรคตับวายเฉียบพลัน (AHPND/EMS)}\index{AHPND (EMS)}):} เกิดจากเชื้อแบคทีเรีย \textit{Vibrio parahaemolyticus} สายพันธุ์ที่มีพลาสมิดสร้างสารพิษไบนารี PirAB ตับและตับอ่อน (Hepatopancreas) ของกุ้งจะมีอาการฝ่อลีบ ซีดขาว ไม่มีอาหารในกระเพาะและลำไส้
    \item \textbf{โรคขี้ขาว (White Feces Syndrome หรือ WFS):} สัมพันธ์กับการติดเชื้อไมโครสปอริเดีย \textit{Enterocytozoon hepatopenaei} (EHP) ร่วมกับแบคทีเรียกลุ่มวิบริโอ กุ้งจะขับถ่ายมูลเป็นเส้นสีขาวลอยตามผิวน้ำ กุ้งโตช้า แคระแกร็น และทยอยตายอย่างต่อเนื่อง
\end{enumerate}

\section{การจำแนกพยาธิสภาพตับกุ้งด้วยโมเดลโครงข่ายเอดจ์}

การนำเทคนิค Deep Learning แบบ \textbf{โครงข่ายประสาทเทียมประสิทธิภาพสูงน้ำหนักเบา (MobileNetV3 / EfficientNet-Lite)} มาพัฒนาเป็นแอปพลิเคชันบนสมาร์ตโฟน ช่วยให้เกษตรกรสามารถถ่ายภาพตับกุ้งและวินิจฉัยโรคได้ทันทีที่ขอบบ่อ

\begin{table}[H]
\centering
\caption{ผลการประเมินสมรรถนะโมเดลจำแนกโรคกุ้ง 4 กลุ่มอาการด้วย Transfer Learning}
\label{tab:shrimp_disease_metrics}
\begin{tabularx}{\textwidth}{p{3.2cm}YYYY}
    \toprule
    \textbf{กลุ่มอาการโรค} & \textbf{ความแม่น (Precision, \%)} & \textbf{การระลึก (Recall, \%)} & \textbf{ค่า F1-Score (\%)} & \textbf{จำนวนภาพทดสอบ ($n$)} \\
    \midrule
    1. ตับปกติสมบูรณ์ (Healthy) & 96.5 & 97.0 & 96.7 & 120 \\
    2. ตัวแดงดวงขาว (WSSV) & 98.2 & 99.0 & 98.6 & 105 \\
    3. ตับฝ่อลีบเฉียบพลัน (AHPND) & 95.8 & 94.5 & 95.1 & 110 \\
    4. ลำไส้ขี้ขาว (EHP/WFS) & 94.0 & 95.2 & 94.6 & 115 \\
    \midrule
    \textbf{ค่าเฉลี่ยรวม (Overall)} & \textbf{96.1} & \textbf{96.4} & \textbf{96.3} & \textbf{450} \\
    \bottomrule
\end{tabularx}
\end{table}

\begin{pythonbox}[การเตรียมโมเดล MobileNetV3 จำแนกพยาธิสภาพกุ้งด้วย PyTorch]
import torch
import torchvision.models as models
import torch.nn as nn

def create_shrimp_classifier(num_classes: int = 4):
    # ใช้สถาปัตยกรรม MobileNetV3-Small โครงสร้างกะทัดรัดสำหรับสมาร์ตโฟน
    model = models.mobilenet_v3_small(pretrained=True)
    
    # ดัดแปลง Linear Classifier ชั้นสุดท้ายให้ตรงกับ 4 คลาสโรคกุ้ง
    in_features = model.classifier[3].in_features
    model.classifier[3] = nn.Sequential(
        nn.Dropout(p=0.2),
        nn.Linear(in_features, num_classes)
    )
    return model

# โมเดลมีขนาดไฟล์เพียง 9.8 MB หลังผ่านการ Quantization เป็น INT8
\end{pythonbox}

\section{การตรวจติดตามสเปกตรัมการดูดกลืนแสงและปรากฏการณ์น้ำเขียว}

สีของน้ำในบ่อเลี้ยงสัตว์น้ำสะท้อนถึงประชากรของแพลงก์ตอนพืช (Phytoplankton) ชนิดที่มีประโยชน์ เช่น ไดอะตอม (Diatoms) จะทำให้น้ำมีสีน้ำตาลอมเขียวหรือสีชาอ่อน ซึ่งเป็นแหล่งอาหารธรรมชาติชั้นเลิศของลูกกุ้ง

ในทางตรงกันข้าม การสะสมของธาตุอาหารไนโตรเจนและฟอสฟอรัสที่มากเกินไปจะกระตุ้นการเพิ่มจำนวนอย่างรวดเร็วของ \textbf{ไซยาโนแบคทีเรียหรือสาหร่ายสีเขียวแกมน้ำเงิน (Cyanobacteria Bloom)} ทำให้น้ำมีสีเขียวเข้มจัด เกิดการสร้างสารพิษไมโครซีสทิน (Microcystin) และมีความเสี่ยงสูงต่อการเกิด \textbf{ปรากฏการณ์น้ำล้มหรือแพลงก์ตอนดรอป (Plankton Crash)\index{ปรากฏการณ์แพลงก์ตอนดรอป (Plankton Crash)}\index{Plankton Crash}} ซึ่งส่งผลให้ออกซิเจนในน้ำดิ่งลงศูนย์เฉียบพลัน

\subsection{การคำนวณดัชนีสเปกตรัมจากภาพถ่ายผิวน้ำ}

การติดตั้งกล้องตรวจวัดสเปกตรัมแสงสะท้อนของผิวน้ำ ช่วยให้สามารถประเมินความเข้มข้นของคลอโรฟิลล์เอ (Chlorophyll-a) โดยการเปรียบเทียบอัตราส่วนแถบแสงสีแดง ($665\text{ nm}$) และแถบแสงอินฟราเรดย่านใกล้ ($705\text{ nm}$)

\begin{equation}
\text{CI}_{\text{green}} = \frac{R_{705}}{R_{665}} - 1
\label{eq:chlorophyll_index}
\end{equation}

เมื่อค่าดัชนี $\text{CI}_{\text{green}}$ พุ่งสูงขึ้นเกินกว่า $3.5$ อย่างรวดเร็วในรอบ 48 ชั่วโมง ระบบจะส่งสัญญาณเตือนภัยการบลูมของสาหร่าย เพื่อให้เกษตรกรลดปริมาณการให้อาหาร เติมจุลินทรีย์บาซิลลัสควบคุม และเตรียมเปิดระบบเติมอากาศล่วงหน้า

\section*{คำถามทบทวนท้ายบทที่ 8}
\addcontentsline{toc}{section}{คำถามทบทวนท้ายบทที่ 8}

\begin{enumerate}
    \item จงอธิบายสาเหตุและลักษณะทางพยาธิสภาพของโรคตัวแดงดวงขาว (WSSV) เทียบกับโรคตับวายเฉียบพลัน (AHPND) ในกุ้งขาวแวนนาไม
    \item เหตุใดสถาปัตยกรรมโครงข่าย MobileNetV3 จึงเหมาะสมกับการติดตั้งใช้งานเพื่อตรวจวินิจฉัยโรคกุ้งบนสมาร์ตโฟนหน้างานจริง
    \item จงอธิบายกลไกการเกิดปรากฏการณ์แพลงก์ตอนดรอป (Plankton Crash) และวิเคราะห์ว่าส่งผลกระทบต่อระดับออกซิเจนในน้ำและของเสียก้นบ่ออย่างไร
    \item การวิเคราะห์สเปกตรัมแสงสะท้อนจากผิวน้ำสามารถช่วยแจ้งเตือนการระบาดของไซยาโนแบคทีเรียล่วงหน้าได้อย่างไร
\end{enumerate}
\clearpage
